% pgfplots figures for the FDP-on-NVMeVirt evaluation.
% Preamble: \usepackage{pgfplots}\pgfplotsset{compat=1.17}
% Reads docs/data/*_summary.csv directly (adjust the relative path if needed).

% --- Over-provisioning sweep: WAF vs working set -----------------------------
\begin{figure}[t]
  \centering
  \begin{tikzpicture}
    \begin{axis}[
        width=.8\linewidth, height=5cm,
        xlabel={Working set (\% of device)}, ylabel={Write amplification (WAF)},
        legend pos=north west, grid=major, mark size=2pt,
    ]
      \addplot table [col sep=comma, x=working_set_pct, y=base_waf] {../data/op_summary.csv};
      \addplot table [col sep=comma, x=working_set_pct, y=fdp_waf]  {../data/op_summary.csv};
      \legend{Baseline, FDP}
    \end{axis}
  \end{tikzpicture}
  \caption{WAF vs.\ over-provisioning. FDP's advantage is largest at moderate
  OP; it vanishes when free space is plentiful (75\%) or nearly exhausted (96\%).}
  \label{fig:fdp-op}
\end{figure}

% --- Lifetime-skew sweep: WAF vs hot-write fraction --------------------------
\begin{figure}[t]
  \centering
  \begin{tikzpicture}
    \begin{axis}[
        width=.8\linewidth, height=5cm,
        xlabel={Hot-write fraction (\%)}, ylabel={Write amplification (WAF)},
        legend pos=north east, grid=major, mark size=2pt,
    ]
      \addplot table [col sep=comma, x=hot_write_pct, y=base_waf] {../data/skew_summary.csv};
      \addplot table [col sep=comma, x=hot_write_pct, y=fdp_waf]  {../data/skew_summary.csv};
      \legend{Baseline, FDP}
    \end{axis}
  \end{tikzpicture}
  \caption{WAF vs.\ workload skew. Baseline WAF is flat; FDP WAF falls as
  lifetimes become more separable, widening the gap.}
  \label{fig:fdp-skew}
\end{figure}
